Classically, Private Information Retrieval (PIR) was studied
in a setting without any pre-processing. In this setting, it is well-known
that 1) public-key cryptography is necessary to achieve   
non-trivial (i.e., sublinear) communication efficiency in the single-server setting, 
and 2) the total server computation 
per query must be linear in the size of the database, no matter in the single-server
or multi-server setting.
Recent works have shown that both of these 
barriers can be overcome if we are willing to introduce a pre-processing phase.
In particular, a recent work called \textsc{Piano} showed
that using only one-way functions, one can construct 
a single-server preprocessing PIR with $\widetilde{O}(\sqrt{n})$
bandwidth and computation per query, 
assuming $\widetilde{O}(\sqrt{n})$ client storage.
For the two-server setting, the state-of-the-art is defined
by two incomparable results. First, \textsc{Piano}
immediately implies a scheme in the two-server setting
with the same performance bounds as stated above. 
Moreover, Beimel et al. showed a 
two-server scheme with $O(n^{1/3})$ bandwidth 
and $O(n/\log^2 n)$ computation per query, and one
with $O(n^{1/2 + \epsilon})$
cost both in bandwidth and computation --- both schemes provide
information theoretic security.


In this paper, we show that assuming the existence of one-way functions,
we can construct a two-server preprocessing PIR scheme 
with $\widetilde{O}(n^{1/4})$ bandwidth and $\widetilde{O}(n^{1/2})$
computation per query,  while requiring only 
$\widetilde{O}(n^{1/2})$ client storage.
We also construct 
a new single-server preprocessing PIR scheme with 
$\widetilde{O}(n^{1/4})$ {\it online} bandwidth and $\widetilde{O}(n^{1/2})$
{\it offline} bandwidth and {\it computation} per query, also requiring 
$\widetilde{O}(n^{1/2})$ client storage.
Specifically, the online bandwidth is the bandwidth required for the client
to obtain an answer, and the offline bandwidth can be viewed as %query-independent 
background maintenance work amortized to each query. 
Our new constructions not only advance the theoretical understanding
of preprocessing PIR, but are also %conceptually simple and 
concretely efficient because the only cryptography
needed is pseudorandom functions. 




%Based on the assumption of the existence of One-Way Function (OWF), we construct a single-server preprocessed PIR scheme with $\tilde{O}(\sqrt{n})$ amortized computation time, $O(\sqrt{n})$ offline communication and $\tilde{O}(n^{1/4})$ online communication per query, while the previous best scheme based on OWF, Piano (S\&P 2024) requires $\Omega(\sqrt{n})$ online communication per query.
%The PIR scheme can also run in the two-server setting and the total amortized communication cost (offline and online) can be improved to $\tilde{O}(n^{1/4})$ per query.

%Also, we design a new technique that allows the state-of-the-art preprocessing-based two-server scheme, TreePIR (CRYPTO 2023), to run in the single-server setting. 
%The resulting PIR scheme is near-optimal, and it has $\tilde{O}(\sqrt{n})$ amortized computation time, $O(\sqrt{n})$ offline communication and polylogarithmic online communication per query, based on the assumption of DDH or LWE.
%Compared to the previous optimal single-server schemes, this scheme does not rely on heavy cryptographic primitives like Privately Punturable Pseudorandom Function or Privately Programmable PRF.

